\documentclass[STT-4300-Main.tex]{subfiles}
\begin{document}
\begin{refsection}

\section{Modèles additifs généralisés (MAG)}
\label{GAM.section}

\subsection{Lissage univarié}
\begin{frame}{Modèles additifs généralisés (MAG)}

On considère le même type de modèles que pour les MLG (distributions de la famille exponentielle pour les variables continues et discrètes, fonctions de lien, etc.), mais on remplace la partie linéaire du MLG par des fonctions non paramétriques :

$$g(\mu_i) = \alpha + f_1(x_{i1}) + \ldots + f_p(x_{ip}).$$

Ce nouveau modèle est particulièrement utile lorsqu'aucune forme fonctionnelle n'est connue a priori. C'est aussi une façon de vérifier si l'hypothèse linéaire (ou une autre hypothèse paramétrique) est raisonnable.

\begin{center}
{\color{blue}Laissons les données nous montrer la forme fonctionnelle appropriée.}
\end{center}
tiré de \citeN{Hast:Tibs:1990}, p.~1.

\end{frame}

\begin{frame}{Lissage univarié (univariate smoothing)}

Données : $(x_i,y_i)$ pour les individus $i=1,\ldots,n$.

\bigskip
On commence par le modèle univarié simple
$$Y_i = f(x_i) + \epsilon_i,$$

avec l'hypothèse que les $\epsilon_i$ sont tirés indépendamment de $N(0, \sigma^2)$.

\bigskip
C'est l'analogue non paramétrique univarié de
$$Y_i = \alpha + \beta x_i + \epsilon_i.$$
\end{frame}


\begin{frame}{Exemple éthanol}

Émissions d'échappement pour 88 combustions d'éthanol dans un moteur d'automobile monocylindre d'essai. Variable réponse : concentration d'oxydes nitriques (\texttt{NOx}). Variable explicative : rapport d'équivalence, qui est une mesure de la richesse du mélange air/éthanol (\texttt{E}).

\begin{center}
\includegraphics[width = 5cm]{figures/ethanol.pdf}  
\end{center}

\end{frame}

\begin{frame}{Ajustement non paramétrique vs quadratique}
\begin{center}
\includegraphics[width = 7cm]{figures/ethanolFit.pdf}  
\end{center}
\end{frame}

\begin{frame}{Fonctions de base (basis functions)}

Comment représenter $f(x)$ ? Avec une approximation à l'aide de fonctions connues :
$$f(x) = \sum_{k=1}^K  b_k(x) \beta_k,$$ où les $b_k(x)$ sont des fonctions de base.
 Le modèle est donc 
 $$Y_i = \sum_{k=1}^K  b_k(x_i) \beta_k + \epsilon_i, $$
 pour $i=1,\ldots,n$.
 \end{frame}

\begin{frame}{Ajustement du modèle}
Pour $Y  = (Y_1, \ldots,Y_n)^T$, $\beta = (\beta_1, \ldots, \beta_K)^T$ et
$\epsilon = (\epsilon_1,\ldots,\epsilon_n)$, le modèle peut être exprimé de façon équivalente comme
$$Y = X \beta + \epsilon,$$
avec
$$X = \left( 
\begin{array}{cccc}
b_1(x_1) & b_2(x_1)  & \ldots  & b_K(x_1) \\
\vdots & \vdots & \vdots & \vdots \\
b_1(x_n) & b_2(x_n)  & \ldots  & b_K(x_n) \\
\end{array} \right).$$

\bigskip
Il peut être ajusté par les moindres carrés :
 $$\hat{\beta} = \mbox{argmin}_{\beta}  ||y - X \beta ||^2.$$
 
\end{frame}

\begin{frame}{Base polynomiale}

Si l'on croit que $f$ est un polynôme d'ordre 4, une base possible est
$$b_1(x) = 1, b_2(x) = x, b_3(x) = x^2, b_4(x) = x^3, b_5(x) = x^4,$$ 
de sorte que 
$$f(x)  = \beta_1 + x \beta_2  + x^2 \beta_3 + x^3 \beta_4 + x^5 \beta_5$$
et le modèle
$$y_i = \beta_1 + x_i \beta_2  + x_i^2 \beta_3 + x_i^3 \beta_4 + x_i^5 \beta_5 + \epsilon_i,$$
\end{frame}

\begin{frame}{La base linéaire par morceaux (ou base tente) -- formule}

Étant donné des nœuds (knots) $x_j^*$, pour $j=1,\ldots,k,$ avec $x_j^* > x_{j-1}^*$, on a pour $j=2, \ldots, k-1$ : 

$$b_j(x) = \left\{ 
\begin{array}{cc}
(x - x_{j-1}^*)/(x_j^* - x_{j-1}^*) & x_{j-1}^* < x < x_j^* \\
 (x_{j+1}^* - x)/(x_{j+1}^* - x_{j}^*) & x_{1}^* < x < x_{j+1}^* \\
 0 & \mbox{sinon}
\end{array} \right.$$

$$b_1(x) = \left\{ 
\begin{array}{cc}
 (x_{2}^* - x)/(x_{2}^* - x_{1}^*) &  x < x_{2}^* \\
 0 & \mbox{sinon}
\end{array} \right.$$

$$ b_k(x) = \left\{ 
\begin{array}{cc}
(x - x_{k-1}^*)/(x_k^* - x_{k-1}^*) &  x > x_{k-1}^* \\
 0 & \mbox{sinon}
\end{array} \right.$$
\end{frame}

\begin{frame}{La base linéaire par morceaux (ou base tente) -- représentation graphique}

\begin{center}
\includegraphics[width = 12cm]{figures/TentBasis.pdf}  
\footnote{Figure tirée des diapositives de S. Wood à l'\textit{Ecole doctorale d'hiver CUSO}, 2014, Les Diablerets}
\end{center}
\end{frame}


\begin{frame}{Approximation de $f$}

Les données $(x_k^* , y_k^*)$ sont interpolées simplement en posant $\beta_k = y_k^*$. La fonction $f$ est représentée en multipliant chaque fonction tente par son coefficient $\beta_k$ et en sommant le résultat.

\begin{center}
\includegraphics[width = 9cm]{figures/TentApprox.pdf}  
\footnote{Figure tirée des diapositives de S. Wood à l'\textit{Ecole doctorale d'hiver CUSO}, 2014, Les Diablerets}
\end{center}
\end{frame}

\begin{frame}{}
En pratique : on choisit certaines valeurs $x_k^*$ réparties sur l'étendue des $x_i$ observés. 

\vspace*{-0.5cm}
\begin{center}
\includegraphics[width = 7.5cm]{figures/EthanolK.pdf}  
\end{center}

\end{frame}


\begin{frame}{Autres bases : les splines}
Les bases polynomiales et tentes sont faciles à comprendre, mais il est possible de faire mieux. En particulier, les bases de splines réduisent l'erreur d'approximation de la fonction pour une dimension donnée de la base de lissage.

\bigskip  
Une spline de lissage cubique (cubic smoothing spline) résout (par rapport à $f$)
$$\sum_{i=1}^n \left(y_i - f(x_i) \right)^2 + \lambda \int (f^{\prime
\prime}(x))^2 dx.$$

Les splines cubiques découlent naturellement de la spécification de la fonction objectif de lissage ci-dessus (définie de manière indépendante de la base).


\end{frame}

\begin{frame}{}
L'inconvénient des splines de lissage (smoothing splines) est qu'elles ont $n$ paramètres libres. Pour conserver leurs bonnes propriétés tout en améliorant l'efficacité computationnelle, on utilise plutôt des splines de régression pénalisées (penalized regression splines).

\bigskip
Cela implique de construire une base de splines pour un jeu de données beaucoup plus petit, puis d'utiliser cette base (avec des pénalités) pour modéliser le jeu de données original. Les valeurs des covariables dans le jeu de données réduit doivent être choisies de manière à bien couvrir l'étendue des covariables du jeu de données original.

\end{frame}



\begin{frame}{Vérification informelle de la dimension de la base}

On effectue une vérification pour s'assurer que $k$ n'est pas trop petit.

\begin{enumerate}
\item Examiner les résidus partiels (partial residuals) par rapport aux estimations des termes, en cherchant des écarts systématiques.
\item Utiliser un test de randomisation des résidus pour détecter une structure résiduelle.\\
Idée : estimer le paramètre d'échelle $\sigma^2$ par différenciation des résidus. La différenciation des résidus qui sont voisins selon $x_i$ devrait donner une estimation de $\sigma^2$ indistinguable d'une estimation par différenciation obtenue avec un ordonnancement aléatoire des résidus, sous l'hypothèse nulle qu'il n'y a pas de structure résiduelle. 
\end{enumerate}
\end{frame}


\begin{frame}{Exemple éthanol avec des splines de régression (regression splines)}

\vspace*{-0.5cm}
\begin{center}
\includegraphics[width = 7.5cm]{figures/EthanolKsplines.pdf}  
\end{center}
\end{frame}

\begin{frame}{Résidus partiels}

\vspace*{-0.5cm}
\begin{center}
\includegraphics[width = 7.5cm]{figures/EthanolKresid.pdf}  
\end{center}
\end{frame}

\begin{frame}[fragile]{}

{\scriptsize
\begin{verbatim}
ethanol1 <- gam(NOx~s(E,k=5),data=ethanol)

> gam.check(ethanol1)

Method: GCV   Optimizer: magic
Smoothing parameter selection converged after 10 iterations.
The RMS GCV score gradient at convergence was 3.154738e-06 .
The Hessian was positive definite.
Model rank =  5 / 5 

Basis dimension (k) checking results. Low p-value (k-index<1) may
indicate that k is too low, especially if edf is close to k'.

                k'   edf k-index p-value
s(ethanol$E) 4.000 3.982   0.682       0
\end{verbatim}
}
\end{frame}


\begin{frame}[fragile]{}

{\scriptsize
\begin{verbatim}
ethanol2 <- gam(NOx~s(E,k=10),data=ethanol)

> gam.check(ethanol2)

Method: GCV   Optimizer: magic
Smoothing parameter selection converged after 8 iterations.
The RMS GCV score gradient at convergence was 2.421344e-06 .
The Hessian was positive definite.
Model rank =  10 / 10 

Basis dimension (k) checking results. Low p-value (k-index<1) may
indicate that k is too low, especially if edf is close to k'.

                k'   edf k-index p-value
s(ethanol$E) 9.000 7.452   0.766       0
\end{verbatim}
}

\end{frame}

\begin{frame}[fragile]{}

{\scriptsize
\begin{verbatim}
ethanol3 <- gam(NOx~s(E,k=15),data=ethanol)

> gam.check(ethanol3)

Method: GCV   Optimizer: magic
Smoothing parameter selection converged after 8 iterations.
The RMS GCV score gradient at convergence was 1.551372e-05 .
The Hessian was positive definite.
Model rank =  15 / 15 

Basis dimension (k) checking results. Low p-value (k-index<1) may
indicate that k is too low, especially if edf is close to k'.

                 k'    edf k-index p-value
s(ethanol$E) 14.000 10.633   0.884    0.11

\end{verbatim}
}
\end{frame}


\begin{frame}[fragile]{}

{\scriptsize
\begin{verbatim}
ethanol4 <- gam(NOx~s(E,k=20),data=ethanol)

> gam.check(ethanol4)

Method: GCV   Optimizer: magic
Smoothing parameter selection converged after 6 iterations.
The RMS GCV score gradient at convergence was 1.246923e-06 .
The Hessian was positive definite.
Model rank =  20 / 20 

Basis dimension (k) checking results. Low p-value (k-index<1) may
indicate that k is too low, especially if edf is close to k'.

                 k'    edf k-index p-value
s(ethanol$E) 19.000 10.603   0.872    0.09
\end{verbatim}
}
\end{frame}

\begin{frame}{Un choix plus formel de $K$ ?}

\begin{itemize}
\item Les modèles avec différentes valeurs de $K$ ne sont pas emboîtés : cela exclut les tests d'hypothèses.
\item Il faut ajuster toutes les valeurs possibles de $K$ si l'on veut utiliser, par exemple, l'AIC.
\item Difficile à généraliser aux modèles avec plus d'une fonction.
\end{itemize}

\bigskip
De plus, où placer les nœuds ?

\end{frame}


\begin{frame}{Lissage (smoothing)}

Au lieu de sélectionner $K$, on utilise plutôt le \textit{lissage} :

\begin{itemize}
\item Choisir $K$ « suffisamment grand » (biais négligeable).
\item Utiliser des $x_k^*$ également espacés.
\item Pour éviter le surajustement, pénaliser l'oscillation (wiggliness) de $f$, par exemple avec 

$${\cal P} (f) = \sum_{j=2}^{k-1} \{f(x_{j-1}^*) - 2 f(x_j^*) + f(x_{j+1}^*) \}^2.$$
Cette pénalité peut être vue comme une approximation grossière de la dérivée seconde de $f$.

\bigskip
Lorsque $f$ est oscillante, la pénalité est grande ; lorsque $f$ est lisse, la pénalité est faible. (La pénalité est nulle pour une droite.)
\end{itemize}
\end{frame}

\begin{frame}{}
Le modèle est ajusté par minimisation du critère des moindres carrés pénalisés
 $$\hat{\beta} = \mbox{argmin}_{\beta} \{  ||y - X \beta ||^2 + \lambda {\cal P}(f) \}.$$

\bigskip
Le \textit{paramètre de lissage} (smoothing parameter) $\lambda$ contrôle le compromis entre la régularité de $\hat{f}$ estimé et la fidélité aux données.

\bigskip
 $\lambda \rightarrow \infty$ mène à une estimation linéaire de $f$, tandis que $\lambda=0$ donne une estimation par régression linéaire par morceaux non pénalisée.
\end{frame}

\begin{frame}{Lissage avec la base tente}
Pour la base tente, on a $\beta_j = f(x_j^*)$ et donc la pénalité peut être exprimée comme une forme quadratique, car
$$ \left( 
\begin{array}{c}
\beta_1 - 2 \beta_2 + \beta_3 \\
\beta_2 - 2 \beta_3 + \beta_4 \\
\beta_3 - 2 \beta_4 + \beta_5 \\
\vdots
\end{array}
\right) = \left[
\begin{array}{ccccccc}
1 & -2 & 1 & 0 & . & . & . \\
0 & 1 & -2 & 1 & 0 & . & . \\
0 & 0 & 1 & -2 & 1 & 0 & . \\
\vdots & \vdots &\vdots &\vdots &\vdots &\vdots &\vdots \\
\end{array}
\right]  \left[
\begin{array}{c}
\beta_1 \\
\beta_2 \\
\beta_3 \\
\vdots
\end{array}
\right] = D \beta
$$
et donc 
$$ {\mathcal P}(f) = \sum_{j=2}^{k-1} \{f(x_{j-1}^*) - 2 f(x_j^*) + f(x_{j+1}^*) \}^2 = \beta^T D^T D \beta = \beta^T S \beta$$
\end{frame}

\begin{frame}{Lissage (général)}
On résout
 $$\hat{\beta} = \mbox{argmin}_{\beta} \{  ||y - X \beta ||^2 + \lambda  \beta^T S \beta \},$$
qui a la solution explicite
 $$\hat{\beta} = (X^T X + \lambda S)^{-1} X^T y.$$
 
\bigskip
Par conséquent 
$$ \hat{y} = \hat{f}(x) = X \hat{\beta} = X (X^T X + \lambda S)^{-1} X^T y = A y,$$
où $A= X (X^T X + \lambda S)^{-1} X^T$ est la matrice d'influence (ou matrice chapeau) (influence/hat matrix). 

\label{linear.smoother}

\end{frame}

\begin{frame}{}
Pour la stabilité numérique, on note que le problème est transformé sous une forme augmentée comme suit :
 $$\{  ||y - X \beta ||^2 + \lambda  \beta^T S \beta \} = \left|\left| \left[
 \begin{array}{c}
y \\
0
\end{array}
 \right] - \left[
 \begin{array}{c}
 X  \\
 \sqrt{\lambda} D
  \end{array}
 \right] \beta \ 
  \right|\right|^2,$$
ce qui est un problème de moindres carrés.
\end{frame}


\begin{frame}{Base naturelle (natural basis)}

Pour toute matrice du modèle $X$ et matrice de pénalité $S$, on peut considérer la reparamétrisation qui rend la matrice de pénalité diagonale.

\bigskip
Soit $X= Q R$ (décomposition QR) et $R^{-T} S R^{-1} = U \Lambda U^{T}$ (décomposition en valeurs singulières).

\bigskip
On définit $P = U^T R$ et on reparamétrise $\beta^{\prime} = P \beta$.

\bigskip 
Dans la nouvelle paramétrisation, la matrice du modèle est $X^{\prime} = Q U$, qui a des colonnes orthogonales ($X = X^{\prime}P$). La matrice de pénalité est maintenant la matrice diagonale $\Lambda$.
\end{frame}


\begin{frame}{Degrés de liberté}
La pénalisation restreint la liberté de variation des coefficients : $K$ coefficients ont $< K$ \textit{degrés de liberté effectifs} (effective degrees of freedom) (DLE).

\bigskip 
On considère une paramétrisation naturelle. 

\bigskip 
Sans pénalisation : $ \tilde{\beta}^{\prime} =  ( X^{\prime} )^T  y$
 
Avec pénalisation : $\hat{\beta}^{\prime} = (I + \lambda \Lambda)^{-1}  ( X^{\prime} )^T
y$, donc $\hat{\beta}_j ^\prime = 
\tilde{\beta}_j^\prime (1 +\lambda \Lambda_{jj})^{-1}$. Ainsi, $(1 + \lambda \Lambda_{jj})^{-1}$ est le facteur de rétrécissement (shrinkage factor) du $j^{\text{e}}$ coefficient. Il donne les DLE pour $\hat{\beta}_j$.

\bigskip
Les DLE totaux sont $\sum_{j} (1 + \lambda  \Lambda_{jj})^{-1} = Tr(F)$ avec $F = (X^T X +
\lambda S)^{-1} X^T X$. 
\end{frame}


\begin{frame}{Biais du lissage}

Puisque $\hat{\beta} = (X^T X + \lambda S)^{-1} X^T y$, on a
\begin{eqnarray*}
E(\hat{\beta}) & = & (X^T X + \lambda S)^{-1} X^T E(y) \\
 & = & (X^T X + \lambda S)^{-1} X^T X \beta \\
& = & F \beta \neq \beta
\end{eqnarray*}

Les lissages sont biaisés en raison de la pénalisation.

\bigskip
Ce biais rend l'inférence fréquentiste difficile.
\end{frame}

\begin{frame}{Le paradigme bayésien en bref}

Avec la distribution $f(y \mid \beta) $ de $(y \mid \beta)$ 
et la {\color{blue}loi a priori (prior distribution)} $\pi(\beta)$ sur le paramètre $\beta$,
on construit une 
{\color{blue}loi a posteriori (posterior distribution)}, c'est-à-dire la 
distribution de $\beta \mid y$, qui est (grâce au théorème de Bayes) :
$$  \pi(\beta \mid y) = \frac{ f(y \mid \beta) \pi(\beta)}{f(y)} .$$


\bigskip
L'inférence est effectuée sur la base de la loi a posteriori.

\end{frame}

\begin{frame}{Illustration du paradigme bayésien}

\begin{center}
\includegraphics[width = 7.5cm]{figures/BayesParadigm.pdf}  
\end{center}

\vspace*{-1cm}
{\tiny Tiré de : 
\texttt{https://cdn1.sph.harvard.edu/wp-content/uploads/sites/565/2018/08/233Spr15\char`_Part1\char`_Bayes.pdf}
}
\end{frame}


\begin{frame}{Modèle de lissage bayésien}
La pénalisation peut être vue comme une loi a priori sur l'oscillation dans un cadre bayésien :

$$\mbox{loi a priori sur l'oscillation} \propto \exp\left( - \lambda \beta^T S \beta /(2
\sigma^2)\right)$$
équivalent à une loi a priori $\beta \sim {\mathcal N}(0,S^{-} \sigma^2/\lambda)$, où
$S^{-} $ est un inverse généralisé de $S$.

\bigskip
Du modèle, on a $y\mid \beta \sim  {\mathcal N}(X \beta, \sigma^2 I)$,
donc (règle de Bayes) 
$$\beta \mid y \sim {\mathcal N}(\hat{\beta}, \sigma^2 (X^T X + \lambda
S)^{-1}).$$

\bigskip
Aussi : $\hat{\sigma} ^2 = ||y - X \hat{\beta}||^2/ (n - Tr(F))$.

\end{frame}

\begin{frame}{Exemple éthanol avec différents $\lambda$ (base tente)}

\vspace*{-0.5cm}
\begin{center}
\includegraphics[width = 7.5cm]{figures/Ethanollambda.pdf}  
\end{center}
\end{frame}

\begin{frame}{Exemple éthanol avec différents $\lambda$ (base spline)}

\vspace*{-0.5cm}
\begin{center}
\includegraphics[width = 7.5cm]{figures/Ethanollambdaspline.pdf}  
\end{center}
\end{frame}


\begin{frame}{Choix de $\lambda$ ?}

Idéalement, on choisit $\lambda$ de sorte que $\hat{f}$ soit aussi proche que possible de $f$, en minimisant, par exemple,

$$M = \frac{1}{n} \sum_{i=1}^n \left( \hat{f}(x_i) - f(x_i) \right)^2.$$

\bigskip

Mais $f$ est inconnue et $M$ ne peut pas être utilisé directement.
\end{frame}

\begin{frame}{Estimation de $M$ par validation croisée (cross-validation)}

On définit le critère de validation croisée ordinaire
$$OCV(\lambda) = \frac{1}{n} \sum_{i=1}^n (y_i - \hat{f}^{-i}_{\lambda}(x_i))^2,$$
où $\hat{f}^{-i}_{\lambda} (x_i)$ désigne l'ajustement en $x_i$ en omettant la $i$-ème observation. Cela implique l'ajustement de $n$ sous-modèles.

\bigskip
On peut démontrer que $E(OCV) \approx E(M) + \sigma^2$.

\end{frame}

\begin{frame}{}

Puisque $\hat{y} = A y$ (voir page~\pageref{linear.smoother})
$$OCV(\lambda) =  \frac{1}{n} \sum_{i=1}^n \Big(\frac{y_i -
  \hat{f}_{\lambda}(x_i)}{1-A_{ii}} \Big)^2,$$
ce qui peut être obtenu à partir de l'ajustement du jeu de données complet uniquement.

\bigskip
En pratique, les $A_{ii}$ sont souvent remplacés par leur moyenne, ce qui donne le critère de validation croisée généralisée (generalized cross-validation) 
$$GCV(\lambda) =  \frac{1}{n} \sum_{i=1}^n \Big(\frac{y_i -
  \hat{f}_{\lambda}(x_i)}{1-Tr(A)/n} \Big)^2 = \frac{n \sum_{i=1}^n \Big(y_i -
  \hat{f}_{\lambda}(x_i)\Big)^2}{(n-Tr(A))^2}.$$

\bigskip
Note : $Tr(A) = Tr(F)$ en raison de la circularité de la trace.
\end{frame}

\begin{frame}{GCV pour le jeu de données éthanol (base spline)}

\vspace*{-0.5cm}
\begin{center}
\includegraphics[width = 7.5cm]{figures/EthanolGCV.pdf}  
\end{center}
\end{frame}


\begin{frame}[fragile]{Ajustement automatique éthanol}

{\scriptsize
\begin{verbatim}
> ethanol.gam <- gam(NOx~s(E,k=15),data=ethanol)

> ethanol.gam
Family: gaussian 
Link function: identity 

Formula:
NOx ~ s(E, k = 15)

Estimated degrees of freedom:
10.6  total = 11.63 

GCV score: 0.1064968 

> ethanol.gam$sp
       s(E) 
0.001490695   
\end{verbatim}
}

\end{frame}

\begin{frame}[fragile]{Ajustement automatique éthanol}

{\tiny
\begin{verbatim}
> summary(ethanol.gam)

Family: gaussian 
Link function: identity 

Formula:
NOx ~ s(E, k = 15)

Parametric coefficients:
            Estimate Std. Error t value Pr(>|t|)    
(Intercept)  1.95737    0.03241    60.4   <2e-16 ***
---
Signif. codes:  0 '***' 0.001 '**' 0.01 '*' 0.05 '.' 0.1 ' ' 1

Approximate significance of smooth terms:
       edf Ref.df     F p-value    
s(E) 10.63  12.31 91.39  <2e-16 ***
---
Signif. codes:  0 '***' 0.001 '**' 0.01 '*' 0.05 '.' 0.1 ' ' 1

R-sq.(adj) =  0.928   Deviance explained = 93.7%
GCV = 0.1065  Scale est. = 0.092419  n = 88
\end{verbatim}
}

\end{frame}

\begin{frame}{Ajustement final}

\vspace*{-0.5cm}
\begin{center}
\includegraphics[width = 7.5cm]{figures/Ethanolfinal.pdf}  
\end{center}
\end{frame}

\subsection{Modèles additifs}
\begin{frame}{Modèles additifs (\textit{additive models}) -- MA}
Comment passe-t-on de 1 à $p$ prédicteurs?

\bigskip
Idéalement, il faudrait considérer des lisseurs de surface
$$y = f(x_1, \ldots, x_p) + \epsilon,$$
mais lorsque $p$ est grand, cet objectif est hors de portée.

\bigskip 
Les {\color{blue}modèles additifs} sont une approximation~:
$$y=\alpha+f_1(x_1) + f_2(x_2)+\ldots+f_p(x_p)+\epsilon,$$
avec $E(\epsilon)=0$, $Var(\epsilon)=\sigma^2$ et $\epsilon$ indépendant de $x_j$.

\end{frame}


\begin{frame}{}

\bigskip
Avec plus d'une composante, il y a un problème d'identifiabilité (\textit{identifiability problem})~: chaque fonction
$f_j$ n'est estimable qu'à une constante additive près. La contrainte
$$\sum_{i=1}^n f_j(x_{ji})=0$$
pour tout $j$ est donc ajoutée. 

\bigskip
Les extensions du modèle additif simple ci-dessus peuvent inclure des fonctions composantes
à deux dimensions ou plus, des termes de variables catégorielles et leurs interactions
avec des variables continues.
\end{frame}

\begin{frame}{}
\textbf{Remarques}~:
\begin{itemize} 
\item La forme additive du modèle permet de tracer les $p$ fonctions
séparément (meilleure interprétabilité).
\item Les modèles additifs sont des approximations plus générales que les modèles de régression linéaire.
\item Les ajustements additifs peuvent être utiles pour tester la linéarité.
\end{itemize}

\end{frame}


\begin{frame}{Ajustement des modèles additifs}

On représente chaque fonction à l'aide de fonctions de base telles que $f_j(x_j) =
\sum_{k=1}^K b_k(x_j) \delta_k^{(j)}$. Le modèle est
$$y_i = \alpha+ \sum_{j=1}^p f_j(x_{ji}) +\epsilon_i = \alpha+ \sum_{j=1}^p X^{(j)} \delta^{(j)} +\epsilon_i,$$
où $X^{(j)}$ contient les fonctions de base pour la covariable $j$. 

\bigskip 
On définit les pénalités associées $(\delta^{(j)})^T \bar{S}_j \delta^{(j)}$ et $X
= [1 \  X^{(1)}  \ldots  \ X^{(p)} ]$. On considère aussi le vecteur de tous les
paramètres $\beta = (\alpha, \delta^{(1)}, \ldots, \delta^{(p)})^T$.
\end{frame}



\begin{frame}{}
On résout maintenant
 $$\hat{\beta} = \mbox{argmin}_{\beta} \{  ||y - X \beta ||^2 + \sum_{j=1}^{p} \lambda_j  \beta^T S_j \beta \},$$
où $S_j$ est construit à partir de $\bar{S}_j$ et complété par des zéros de sorte que
$(\delta^{(j)})^T \bar{S}_j \delta^{(j)} = \beta^T S_j \beta$.

\bigskip
De façon similaire au cas univarié, on a
$$\hat{\beta} = (X^T X + \lambda_1 S_1 + \ldots + \lambda_p S_p)^{-1} X^T y,$$
$$A = X (X^T X + \lambda_1 S_1 + \ldots + \lambda_p S_p)^{-1} X^T,$$
et
$$F = (X^T X + \lambda_1 S_1 + \ldots + \lambda_p S_p)^{-1} X^T X .$$
\end{frame}


\begin{frame}{}
Pour la stabilité des calculs, le problème augmenté correspondant est considéré~:
\begin{eqnarray*}
& & \left|\left| y - X \beta \right|\right|^2 + \beta^T (\lambda_1 S_1 +
\ldots + \lambda_p S_p) \beta \\
& & = \left|\left|
\left(\begin{array}{c}
y \\
0
\end{array} \right) -
\left(\begin{array}{c}
X \\
B
\end{array} \right) \beta
\right|\right|^2,
\end{eqnarray*}
où $B$ est tel que $B^T B = \lambda_1 S_1 + \ldots + \lambda_p S_p.$

\bigskip
C'est un problème de moindres carrés.
\end{frame}

\begin{frame}{Choix de $(\lambda_1, \ldots, \lambda_p)$ par GCV}

Trouver $(\lambda_1, \ldots, \lambda_p)$ qui minimise
\begin{eqnarray}
GCV(\lambda_1, \ldots, \lambda_p) & = &   \frac{1}{n} \sum_{i=1}^n \Big(\frac{y_i -
  \hat{f}_{\lambda_1, \ldots, \lambda_p}(x_i)}{1-Tr(A)/n} \Big)^2 \\
  & = &  \frac{n \sum_{i=1}^n \Big(y_i -
  \hat{f}_{\lambda_1, \ldots, \lambda_p}(x_i)\Big)^2}{(n-Tr(A))^2}.
\end{eqnarray}

\end{frame}



\begin{frame}{Distribution a posteriori de $\beta$}
On a
$$ \beta | y \sim {\mathcal N}(\hat{\beta}, V_{\beta}),$$
où $V_{\beta} = \sigma^2 (X^T X + \lambda_1 S_1 + \ldots + \lambda_p S_p)$.

\bigskip 
$V_{\beta}$ est estimé en substituant $\hat{\sigma}^2 = ||y - X
\hat{\beta}||^2/ (n - Tr(F))$. 

\bigskip
Le résultat ci-dessus peut être utilisé pour l'inférence sur $\beta$, soit directement (quantités linéaires en les paramètres du modèle), soit par simulation ou bootstrap (dans les autres cas).
\end{frame}

\begin{frame}{Tests d'hypothèses sur $\beta$}
Supposons qu'on veuille tester $H_0: \beta_j=0$, où $\beta_j$ est un sous-vecteur de taille
$p_j$ de 
$\beta$. Soit $V_{\beta_j}$ le bloc de $V_{\beta}$ correspondant à
$\beta_j$. 

\bigskip
Alors, sous $H_0$, on peut utiliser la statistique de test
$$\hat{\beta_j}^T V_{\beta_j}^{-1} \hat{\beta_j}/p_j \sim F_{p_j,n-p},$$
s'il y a un paramètre d'échelle ($\phi$ inconnu), et 
$$\hat{\beta_j}^T V_{\beta_j}^{-1} \hat{\beta_j} \sim \chi_{p_j}^2,$$
sinon ($\phi$ connu).
\end{frame}

\begin{frame}{Intervalles de confiance pour $f_j$}

À partir de la distribution a posteriori de $\beta$, des intervalles de crédibilité peuvent être construits pour $f_j$. 

\bigskip
Soit $\tilde{f}_j = \tilde{X} \beta$, où $\tilde{X}$ a des zéros dans les colonnes correspondant aux coefficients n'ayant rien à voir avec $f_j$, tandis que ses autres colonnes contiennent les fonctions de base pour $f_j$. Soit $v= \mbox{diag}(\tilde{X} V_{\beta} \tilde{X}^T)$, alors
$$\hat{f_{ji}} \pm z_{\alpha/2} \sqrt{v_i}$$
est un intervalle de crédibilité approximatif à $(1-\alpha) 100 \%$ pour $f_{ji}$.

\bigskip
Il s'avère que les intervalles de crédibilité bayésiens possèdent de bonnes propriétés de couverture fréquentistes.

\end{frame}


\begin{frame}{Sélection de variables (\textit{variable selection}) pour les termes lisses}
Une statistique de test peut être dérivée pour tester l'hypothèse nulle $H_0:
f_j(x_j)=0$. Soit $f_j = \tilde{X} \beta$. Alors, de $\hat{\beta} \sim {\mathcal
N}(\beta, V_{\beta})$, on a $ \hat{f_j} \sim {\mathcal N}(f_j,  V_{f_j}),$
où $V_{f_j} = \tilde{X}  V_{\beta} \tilde{X}^T$.

\bigskip
La statistique de test est définie par
$$ T_r = \hat{f_j}^T V_{f_j}^{r -} \hat{f_j} ,$$
où $V_{f_j}^{r -}$ est une pseudo-inverse de rang $r$ de $V_{f_j}$.

\bigskip
Il faut prendre certaines précautions pour le choix de $r$. En bref, ce choix se fait en fonction des degrés de liberté estimés de la fonction $f_j$.

\bigskip
La distribution de $T_r$ sous $H_0$ est $\chi^2_r$ si $r$ est entier.
\end{frame}

\begin{frame}[fragile]{Analyse des données Boston housing}

{\tiny
\begin{verbatim}
> BHnpfit <- gam(log(medv)~chas+s(crim)+s(zn)+s(indus)+s(nox)+s(rm)+s(age)+s(dis)+
   s(rad,k=8)+s(tax)+s(ptratio)+s(b)+s(lstat),data=BostonHousing)
> BHnpfit

Family: gaussian 
Link function: identity 

Formula:
log(medv) ~ chas + s(crim) + s(zn) + s(indus) + s(nox) + s(rm) + 
    s(age) + s(dis) + s(rad, k = 8) + s(tax) + s(ptratio) + s(b) + 
    s(lstat)

Estimated degrees of freedom:
3.12 1.00 7.83 9.00 4.59 1.00 8.61 
1.33 3.09 1.00 4.99 5.24  total = 52.8 

GCV score: 0.02158189     
\end{verbatim}
}

\end{frame}

\begin{frame}[fragile]{}
{\tiny
\begin{verbatim}
> gam.check(BHnpfit)

Method: GCV   Optimizer: magic
Smoothing parameter selection converged after 20 iterations.
The RMS GCV score gradient at convergence was 2.97197e-08 .
The Hessian was positive definite.
Model rank =  108 / 108 

Basis dimension (k) checking results. Low p-value (k-index<1) may
indicate that k is too low, especially if edf is close to k'.

              k'   edf k-index p-value
s(crim)    9.000 3.123   1.074    0.94
s(zn)      9.000 1.000   0.797    0.00
s(indus)   9.000 7.825   0.790    0.00
s(nox)     9.000 9.000   0.838    0.00
s(rm)      9.000 4.591   0.941    0.10
s(age)     9.000 1.000   0.974    0.23
s(dis)     9.000 8.613   0.936    0.08
s(rad)     7.000 1.330   0.796    0.00
s(tax)     9.000 3.089   0.798    0.00
s(ptratio) 9.000 1.000   0.814    0.00
s(b)       9.000 4.989   0.981    0.34
s(lstat)   9.000 5.240   1.072    0.94
\end{verbatim}
}
\end{frame}

\begin{frame}[fragile]{}

{\tiny
\begin{verbatim}
> BHnpfit.new <- gam(log(medv)~chas+s(crim)+s(zn)+s(indus)+s(nox,k=20)+s(rm)+s(age)+
    s(dis,k=20)+s(rad,k=8)+s(tax)+s(ptratio)+s(b)+s(lstat),data=BostonHousing)
> BHnpfit.new

Family: gaussian 
Link function: identity 

Formula:
log(medv) ~ chas + s(crim) + s(zn) + s(indus) + s(nox, k = 20) + 
    s(rm) + s(age) + s(dis, k = 20) + s(rad, k = 8) + s(tax) + 
    s(ptratio) + s(b) + s(lstat)

Estimated degrees of freedom:
 3.37  1.00  5.21 13.07  5.10  1.00 13.40 
 2.26  3.20  1.00  1.94  4.76  total = 57.3 

GCV score: 0.02116857 
\end{verbatim}
}
\end{frame}
 

\begin{frame}[fragile]{}

{\tiny
\begin{verbatim}
> gam.check(BHnpfit.new)

Method: GCV   Optimizer: magic
Smoothing parameter selection converged after 21 iterations.
The RMS GCV score gradient at convergence was 2.186018e-08 .
The Hessian was positive definite.
Model rank =  128 / 128 

Basis dimension (k) checking results. Low p-value (k-index<1) may
indicate that k is too low, especially if edf is close to k'.

               k'    edf k-index p-value
s(crim)     9.000  3.368   1.078    0.94
s(zn)       9.000  1.000   0.838    0.00
s(indus)    9.000  5.211   0.819    0.00
s(nox)     19.000 13.065   0.875    0.00
s(rm)       9.000  5.101   0.948    0.14
s(age)      9.000  1.000   0.990    0.40
s(dis)     19.000 13.397   0.967    0.20
s(rad)      7.000  2.260   0.835    0.00
s(tax)      9.000  3.202   0.834    0.00
s(ptratio)  9.000  1.000   0.856    0.00
s(b)        9.000  1.945   0.978    0.34
s(lstat)    9.000  4.755   1.077    0.94
\end{verbatim}
}
\end{frame}

\begin{frame}{}

\vspace*{-0.5cm}
\begin{center}
\includegraphics[width = 7.5cm]{figures/BHnpfitnewCheck.pdf}  
\end{center}
\end{frame}

\begin{frame}[fragile]{}

{\tiny
\begin{verbatim}
Family: gaussian 
Link function: identity 
Formula: log(medv) ~ chas + s(crim) + s(zn) + s(indus) + s(nox, k = 20) +  s(rm) + s(age) + 
  s(dis, k = 20) + s(rad, k = 8) + s(tax) +  s(ptratio) + s(b) + s(lstat)

Parametric coefficients:
            Estimate Std. Error t value Pr(>|t|)    
(Intercept) 3.031791   0.006376 475.497   <2e-16 ***
chas        0.039350   0.027264   1.443     0.15    
---
Signif. codes:  0 '***' 0.001 '**' 0.01 '*' 0.05 '.' 0.1 ' ' 1

Approximate significance of smooth terms:
              edf Ref.df      F  p-value    
s(crim)     3.368  4.184 19.183 7.60e-15 ***
s(zn)       1.000  1.000  0.639  0.42437    
s(indus)    5.211  6.119  1.882  0.08638 .  
s(nox)     13.065 15.171 11.476  < 2e-16 ***
s(rm)       5.101  6.262 17.018  < 2e-16 ***
s(age)      1.000  1.000  1.964  0.16174    
s(dis)     13.397 15.676  3.900 9.16e-07 ***
s(rad)      2.260  2.627 10.931 4.05e-06 ***
s(tax)      3.202  3.789 11.648 1.24e-08 ***
s(ptratio)  1.000  1.000 37.846 1.63e-09 ***
s(b)        1.945  2.401  5.254  0.00327 ** 
s(lstat)    4.755  5.864 35.015  < 2e-16 ***
---
Signif. codes:  0 '***' 0.001 '**' 0.01 '*' 0.05 '.' 0.1 ' ' 1

R-sq.(adj) =  0.888   Deviance explained =   90%
GCV = 0.021169  Scale est. = 0.018771  n = 506
\end{verbatim}
}

\end{frame}

\begin{frame}{}

\vspace*{-0.5cm}
\begin{center}
\includegraphics[width = 8cm]{figures/BHnpfitnew.pdf}  
\end{center}
\end{frame}

\begin{frame}[fragile]{Retirer \texttt{zn}, \texttt{age} et \texttt{chas}?}

{\tiny
\begin{verbatim}
> BHnpfit2.new <- gam(log(medv)~s(crim)+s(indus)+s(nox,k=20)+s(rm)+s(dis,k=20)+
   s(rad,k=8)+s(tax)+s(ptratio)+s(b)+s(lstat),data=BostonHousing)
> BHnpfit2.new

Family: gaussian 
Link function: identity 

Formula: log(medv) ~ s(crim) + s(indus) + s(nox, k = 20) + s(rm) + s(dis, k = 20) + 
   s(rad, k = 8) + s(tax) + s(ptratio) + s(b) + s(lstat)

Estimated degrees of freedom:
 3.35  5.12 12.79  4.96 13.33  2.35  3.27 
 1.00  1.91  4.83  total = 53.92 

GCV score: 0.02112012 
\end{verbatim}
}

\end{frame}


\begin{frame}[fragile]{}

{\tiny
\begin{verbatim}
> gam.check(BHnpfit2.new)

Method: GCV   Optimizer: magic
Smoothing parameter selection converged after 20 iterations.
The RMS GCV score gradient at convergence was 4.284208e-08 .
The Hessian was positive definite.
Model rank =  109 / 109 

Basis dimension (k) checking results. Low p-value (k-index<1) may
indicate that k is too low, especially if edf is close to k'.

               k'    edf k-index p-value
s(crim)     9.000  3.352   1.082    0.96
s(indus)    9.000  5.123   0.813    0.00
s(nox)     19.000 12.795   0.876    0.00
s(rm)       9.000  4.961   0.940    0.08
s(dis)     19.000 13.325   0.967    0.22
s(rad)      7.000  2.353   0.828    0.00
s(tax)      9.000  3.268   0.823    0.00
s(ptratio)  9.000  1.000   0.849    0.00
s(b)        9.000  1.915   0.987    0.44
s(lstat)    9.000  4.828   1.076    0.96
\end{verbatim}
}

\end{frame}

\begin{frame}{}

\vspace*{-0.5cm}
\begin{center}
\includegraphics[width = 7.5cm]{figures/BHnpfitnew2Check.pdf}  
\end{center}
\end{frame}

\begin{frame}[fragile]{}

{\tiny
\begin{verbatim}
> summary(BHnpfit2.new)

Family: gaussian 
Link function: identity 

Formula: log(medv) ~ s(crim) + s(indus) + s(nox, k = 20) + s(rm) + s(dis, k = 20) + 
   s(rad, k = 8) + s(tax) + s(ptratio) + s(b) + s(lstat)

Parametric coefficients:
            Estimate Std. Error t value Pr(>|t|)    
(Intercept) 3.034513   0.006107   496.9   <2e-16 ***
---
Signif. codes:  0 '***' 0.001 '**' 0.01 '*' 0.05 '.' 0.1 ' ' 1

Approximate significance of smooth terms:
              edf Ref.df      F  p-value    
s(crim)     3.352  4.165 19.885 2.34e-15 ***
s(indus)    5.123  6.030  2.132  0.04962 *  
s(nox)     12.795 14.900 11.670  < 2e-16 ***
s(rm)       4.961  6.109 17.497  < 2e-16 ***
s(dis)     13.325 15.611  3.932 6.55e-07 ***
s(rad)      2.353  2.735 11.034 2.13e-06 ***
s(tax)      3.268  3.868 11.418 1.14e-08 ***
s(ptratio)  1.000  1.000 44.876 5.87e-11 ***
s(b)        1.915  2.365  5.030  0.00437 ** 
s(lstat)    4.828  5.938 43.788  < 2e-16 ***
---
Signif. codes:  0 '***' 0.001 '**' 0.01 '*' 0.05 '.' 0.1 ' ' 1

R-sq.(adj) =  0.887   Deviance explained = 89.9%
GCV = 0.02112  Scale est. = 0.01887   n = 506
\end{verbatim}
}
\end{frame}

\begin{frame}{}

\vspace*{-0.5cm}
\begin{center}
\includegraphics[width = 8cm]{figures/BHnpfit2new.pdf}  
\end{center}
\end{frame}


\begin{frame}[fragile]{Retirer \texttt{ptratio}?}

{\tiny
\begin{verbatim}
> BHnpfit3.new <- gam(log(medv)~s(crim)+s(indus)+s(nox,k=20)+s(rm)+s(dis,k=20)+s(rad,k=8)
+s(tax)+s(b)+s(lstat),data=BostonHousing)
> BHnpfit3.new

Family: gaussian 
Link function: identity 

Formula:
log(medv) ~ s(crim) + s(indus) + s(nox, k = 20) + s(rm) + s(dis, 
    k = 20) + s(rad, k = 8) + s(tax) + s(b) + s(lstat)

Estimated degrees of freedom:
 3.35  8.10 13.58  5.30 12.33  1.00  3.27 
 1.77  4.70  total = 54.39 

GCV score: 0.02258293   

> BHnpfit3.new$aic
[1] -486.3607
\end{verbatim}
}

La valeur d'Akaike de ce modèle est $-486{,}36$, bien pire que ce qu'on avait auparavant.
\end{frame}


\subsection{Modèles additifs généralisés}

\begin{frame}{Ingrédients des GAM}

\begin{enumerate}
\item Des variables réponse indépendantes $Y_1, \ldots,Y_n$ supposées partager
  la même distribution de la {\color{blue}famille exponentielle} ($F_{\theta_i}$).
\item Un ensemble de {\color{blue}variables explicatives}\\
 $x_i^T = (x_{i1},  \ldots, x_{ip})$ pour $i=1,\ldots,n$.
\item Une fonction de {\color{blue}lien} monotone $g$ telle que
$$g(\mu_i) = \alpha + f_1(x_{i1}) + \ldots + f_p(x_{ip}) = \eta_i.$$
$\eta_i $ est appelé le {\color{blue}prédicteur additif} (\textit{additive predictor}).
\end{enumerate}

\bigskip
On suppose que
$$E(Y_i) = \mu_i \ \mbox{et} \ \ Var(Y_i) = v_i = v(\mu_i).$$
La différence avec le modèle GLM est le fait que $\mu_i$ est obtenu
de manière non paramétrique, via $g^{-1}(\alpha + f_1(x_{i1}) + \ldots + f_p(x_{ip}))$.
\end{frame}

\begin{frame}{}
On a maintenant
$$g(\mu_i ) =  \alpha+ \sum_{j=1}^p f_j(x_{ji}) = \alpha+ \sum_{j=1}^p X^{(j)} \delta^{(j)},$$
avec les pénalités associées $(\delta^{(j)})^T \bar{S}_j \delta^{(j)}$ et $X
= [1 \  X^{(1)}  \ldots  \ X^{(p)} ]$. 

\bigskip 
Considérons $\beta = (\alpha, \delta^{(1)}, \ldots, \delta^{(p)})$. On résout
 $$\hat{\beta} = \mbox{argmin}_{\beta} \{  D(y,\mu)+ \sum_{j=1}^{p} \lambda_j  \beta^T S_j \beta \},$$
où $D(y,\mu)$ est la fonction de déviance (voir p.~\pageref{deviance.def}) et où $S_j$ est construit à partir de $\bar{S}_j$ et complété par des zéros de sorte que
$(\delta^{(j)})^T \bar{S}_j \delta^{(j)} = \beta^T S_j \beta$.
\end{frame}

\begin{frame}{}
Les GAM sont ajustés à l'aide d'un {\color{blue}algorithme de score local} (\textit{local scoring algorithm}).

\bigskip
À l'étape $(m-1)$, on construit une {\color{blue}variable dépendante ajustée} (\textit{adjusted dependent variable})
$$z_i^{(m-1)} = \eta_i^{(m-1)} + (y_i - \mu_i^{(m-1)}) \Big(\frac{\partial \eta_i}{\partial \mu_i} \Big)^{(m-1)},$$
avec $\eta_i^{(m-1)} = \alpha^{(m-1)} + \sum_{j=1}^p f_j^{(m-1)}(x_{ij})$, et des poids
$$w_i^{(m-1)} = \frac{1}{Var(Y_i)^{(m-1)}}\Big(\frac{\partial \mu_i}{\partial \eta_i} \Big)_{(m-1)}^2.$$

On ajuste un {\color{blue}modèle additif pondéré} (\textit{weighted additive model}) à $z_i^{(m-1)}$ pour obtenir les fonctions estimées $f_j^{(m)}$, le prédicteur additif $\eta_i^{(m)}$ et les valeurs ajustées $\mu_i^{(m)}$. On répète jusqu'à convergence.


\bigskip
L'ajustement par moindres carrés pondérés du GLM dans l'algorithme de score local
est remplacé par un ajustement additif pondéré (voir page~\pageref{IWLSweights}).

\end{frame}


\begin{frame}{}
L'ajustement additif pondéré peut s'écrire comme un ajustement par moindres carrés pondérés pour le problème augmenté comme suit~:
$$\left|\left|
\left(\begin{array}{cc}
\sqrt{W^{(m-1)}} & 0 \\
0 & I 
\end{array} \right) 
\left[
\left(\begin{array}{c}
z^{(m-1)} \\
0
\end{array} \right) -
\left(\begin{array}{c}
X \\
B
\end{array} \right) \right]\beta
\right|\right|^2,
$$
avec $W^{(m-1)} = \mbox{diag}(w_1^{(m-1)}, \ldots,w_n^{(m-1)})$ et $B$ tel que $B^T B = \lambda_1 S_1 + \ldots + \lambda_p S_p.$

\end{frame}

\begin{frame}{Propriétés}

La matrice des degrés de liberté effectifs devient
$$F = (X^T  WX + \lambda_1 S_1 + \ldots + \lambda_p S_p)^{-1} X^T W X .$$

\bigskip
On a asymptotiquement
$$ \beta | y \sim {\mathcal N}(\hat{\beta}, V_{\beta}),$$
où $V_{\beta} = \sigma^2 (X^T W X + \lambda_1 S_1 + \ldots + \lambda_p S_p)$.

\bigskip
Le paramètre d'échelle, s'il est inconnu, est estimé par
$\hat{\phi} = \sum_{i=1}^n (y_i - X_i \hat{\beta})^2/ (n - Tr(F))$. 

\end{frame}


\begin{frame}{Choix de $\lambda_1,\ldots,\lambda_p$ pour $\phi$ connu}

Minimiser l'erreur quadratique moyenne~:

\bigskip
Pour un modèle additif à variance constante, on a
\begin{eqnarray*}
M(\lambda_1,\ldots,\lambda_p) & = &  E \left( \frac{ \left|\left| \mu - X \hat{\beta}  \right|\right|^2}{n} \right) \\
& = & E \left( \frac{ \left|\left| y - A y  \right|\right|^2}{n} \right) - \phi + \frac{2 \phi Tr(A)}{n},  
\end{eqnarray*}
(voir p.~52 et p.~255 dans \citeN{Wood:2017} pour l'obtention de cette expression.)

\end{frame}


\begin{frame}{}
La forme de $M$ suggère l'estimateur suivant (estimateur de risque sans biais, \textit{Un-Biased Risk Estimator})
$$ UBRE(\lambda_1,\ldots,\lambda_p) = \frac{ \left|\left| y - Ay  \right|\right|^2}{n}   - \phi + \frac{2 \phi Tr(A)}{n} = D(y,\hat{\mu}) + \frac{2 \phi Tr(F)}{n}.$$
aussi connu sous le nom de $C_p$ de Mallows.

\bigskip 
Le UBRE doit être minimisé.

\bigskip
Par extension, il est utilisé pour toutes les familles avec $\phi$ connu.
\end{frame}

\begin{frame}{Choix de $\lambda$ pour $\phi$ inconnu}

Si $\phi$ est inconnu, il doit être estimé. Substituer $\hat{\phi}$ dans la définition du UBRE le rend inapproprié pour la sélection de modèle.

\bigskip
On utilise donc le $GCV$. Il se généralise à 
$$ GCV(\lambda_1, \ldots,\lambda_p) = \frac{n D(y,\hat{\mu})}{(n- Tr(F))^2}.$$
\end{frame}

\begin{frame}{Validation du modèle}

\begin{itemize}
\item Vérifications informelles de $K$  
\item[]
\item Analyse des résidus
\end{itemize}

\end{frame}

\begin{frame}{Comparaison de modèles (emboîtés)}

Étant dans un cadre de vraisemblance, le critère d'Akaike peut être utilisé pour comparer des modèles emboîtés~:
$$AIC = -2 l(\hat{\beta};y) + 2 \phi Tr(F).$$
Les modèles avec le plus petit $AIC$ sont à privilégier.

\bigskip
$GCV(\hat{\lambda}_1, \ldots,\hat{\lambda}_p)$ estime l'erreur de
prédiction. Il peut être utilisé pour comparer des modèles. Les modèles avec le plus petit $GCV$ sont à privilégier.

\bigskip
La proportion de déviance expliquée est aussi une mesure de la qualité d'un ajustement.

\end{frame}

\subsection{Analyse du jeu de données sur les oiseaux}
\begin{frame}{Analyse du jeu de données sur les oiseaux}

Variable réponse~: nombre d'oiseaux observés (comptages, \texttt{NUMBER}).

\bigskip
Prédicteurs~:
\begin{itemize}
\item année (\texttt{YEAR})  - {\color{blue}Intérêt principal}.
\item jour (\texttt{DAY}) 
\item nombre
d'observateurs sur l'île (\texttt{NO.OBS}) 
\item direction du vent à minuit la nuit précédente (\texttt{WIND.0}) 
\item direction du vent à midi la veille (\texttt{WIND.1}) 
\item direction du vent à minuit 2 jours avant (\texttt{WIND.2}) 
\item direction du vent à midi 2 jours avant (\texttt{WIND.3}) 
\item vitesse du vent ce jour-là (\texttt{WINDSPD}) 
\item condition du ciel (1=dégagé, 2=nuageux, 3=brouillard, 4=pluie) (\texttt{SKY})
\item proportion de lune illuminée (\texttt{MOON}) 
\item espèce d'oiseau (1-5) (\texttt{BIRD})
\end{itemize}

\end{frame}


\begin{frame}[fragile]{Modèle de Poisson avec lien logarithmique}

\begin{eqnarray*}
\lefteqn{\texttt{NUMBER} \sim  \alpha + f_1(\texttt{YEAR}) + f_2(\texttt{DAY}) +
f_3(\texttt{NO.OBS})} \\
& + & f_4(\texttt{WIND.0}) + f_5(\texttt{WIND.1}) +
f_6(\texttt{WIND.2}) \\
& + & f_7(\texttt{WIND.3}) + f_8(\texttt{WINDSPD}) +
\mbox{factor(\texttt{SKY})}\\
& + & f_9(\texttt{MOON}) + \mbox{factor}(\texttt{BIRD})
\end{eqnarray*}

\bigskip

{\scriptsize
\begin{verbatim}
> require(mgcv)
> GrA.fit <- gam(NUMBER~s(YEAR,k=20)+s(DAY,k=20)+s(NO.OBS,k=8)+
   s(WIND.0,k=20)+s(WIND.1,k=20)+s(WIND.2,k=20)+s(WIND.3,k=20)+
   s(WINDSPD,k=20)+factor(SKY)+s(MOON,k=20)+factor(BIRD), 
   family=poisson,data=GrA,na.action=na.omit)
> plot(GrA.fit,scale=0,pages=1,pers=T,cex.lab=2)
\end{verbatim}
}
\end{frame}

\begin{frame}{}

\begin{center}
\includegraphics[width = 8cm]{figures/GrA.pdf}  
\end{center}

\end{frame}

\begin{frame}[fragile]{}

{\scriptsize
\begin{verbatim}
> gam.check(GrA.fit)

Method: UBRE   Optimizer: outer newton
full convergence after 12 iterations.
Gradient range [-1.543031e-05,2.132528e-05]
(score 38.92044 & scale 1).
eigenvalue range [-1.317038e-05,8.758125e-05].
Model rank =  167 / 167 

Basis dimension (k) checking results. Low p-value (k-index<1) may
indicate that k is too low, especially if edf is close to k'.

               k'    edf k-index p-value
s(YEAR)    19.000 18.987   0.676    0.00
s(DAY)     19.000 18.923   0.838    0.00
s(NO.OBS)   7.000  6.941   0.704    0.00
s(WIND.0)  19.000 18.985   0.846    0.00
s(WIND.1)  19.000 18.962   0.894    0.00
s(WIND.2)  19.000 18.951   0.903    0.00
s(WIND.3)  19.000 18.891   0.895    0.00
s(WINDSPD) 19.000 18.982   0.888    0.00
s(MOON)    19.000 18.945   0.959    0.31
\end{verbatim}
}
\end{frame}

\begin{frame}
\begin{center}
\includegraphics[width = 7.5cm]{figures/GrAcheck.pdf}  
\end{center}
\end{frame}

\begin{frame}[fragile]{}

{\scriptsize
\begin{verbatim}
> summary(GrA.fit)

Family: poisson 
Link function: log 

Formula:
NUMBER ~ s(YEAR, k = 20) + s(DAY, k = 20) + s(NO.OBS, k = 8) + 
    s(WIND.0, k = 20) + s(WIND.1, k = 20) + s(WIND.2, k = 20) + 
    s(WIND.3, k = 20) + s(WINDSPD, k = 20) + factor(SKY) + s(MOON, 
    k = 20) + factor(BIRD)

Parametric coefficients:
                  Estimate Std. Error z value Pr(>|z|)    
(Intercept)        2.84578    0.01186 239.898  < 2e-16 ***
factor(SKY)cloudy  0.03137    0.01300   2.413   0.0158 *  
factor(SKY)fog    -0.04209    0.01852  -2.272   0.0231 *  
factor(SKY)rain   -0.09024    0.02229  -4.048 5.16e-05 ***
factor(BIRD)2     -2.00090    0.02604 -76.849  < 2e-16 ***
factor(BIRD)3     -2.60308    0.03415 -76.223  < 2e-16 ***
factor(BIRD)4      0.20671    0.01208  17.118  < 2e-16 ***
factor(BIRD)5      1.22781    0.01020 120.359  < 2e-16 ***
---
Signif. codes:  0 '***' 0.001 '**' 0.01 '*' 0.05 '.' 0.1 ' ' 1

\end{verbatim}
}
\end{frame}

\begin{frame}[fragile]{}

{\scriptsize
\begin{verbatim}
Approximate significance of smooth terms:
              edf Ref.df Chi.sq p-value    
s(YEAR)    18.987 19.000   9207  <2e-16 ***
s(DAY)     18.923 18.999   6365  <2e-16 ***
s(NO.OBS)   6.941  6.998   1862  <2e-16 ***
s(WIND.0)  18.985 19.000   3543  <2e-16 ***
s(WIND.1)  18.962 19.000   2172  <2e-16 ***
s(WIND.2)  18.951 19.000   3784  <2e-16 ***
s(WIND.3)  18.891 18.998   1383  <2e-16 ***
s(WINDSPD) 18.982 19.000   3259  <2e-16 ***
s(MOON)    18.945 18.999   1746  <2e-16 ***
---
Signif. codes:  0 '***' 0.001 '**' 0.01 '*' 0.05 '.' 0.1 ' ' 1

R-sq.(adj) =  0.0895   Deviance explained = 65.5%
UBRE =  38.92  Scale est. = 1         n = 2230
\end{verbatim}
}
\end{frame}


\begin{frame}[fragile]{Modèle quasi-Poisson, base cyclique pour \texttt{WIND}}

{\scriptsize
\begin{verbatim}
> gam.check(GrA.quasifit)
Method: GCV   Optimizer: outer newton
full convergence after 9 iterations.
Gradient range [-3.33374e-06,2.451736e-05]
(score 45.78461 & scale 58.39889).
Hessian positive definite, eigenvalue range [0.003653669,0.04407047].
Model rank =  163 / 163 

Basis dimension (k) checking results. Low p-value (k-index<1) may
indicate that k is too low, especially if edf is close to k'.
               k'    edf k-index p-value
s(YEAR)    19.000 18.819   0.685    0.00
s(DAY)     19.000 12.869   0.831    0.00
s(NO.OBS)   7.000  2.317   0.705    0.00
s(WIND.0)  18.000  9.992   0.840    0.00
s(WIND.1)  18.000  9.003   0.890    0.00
s(WIND.2)  18.000 14.793   0.908    0.00
s(WIND.3)  18.000  4.803   0.887    0.00
s(WINDSPD) 19.000 17.248   0.905    0.00
s(MOON)    19.000 16.376   0.956    0.27
\end{verbatim}
}
\end{frame}

\begin{frame}
\begin{center}
\includegraphics[width = 7.5cm]{figures/GrAquasicheck.pdf}  
\end{center}
\end{frame}

\begin{frame}[fragile]{}

{\scriptsize
\begin{verbatim}
> summary(GrA.quasifit)

Family: quasipoisson 
Link function: log 

Formula:
NUMBER ~ s(YEAR, k = 30) + s(DAY, k = 30) + s(NO.OBS, k = 8) + 
    s(WIND.0, k = 30, bs = "cc") + s(WIND.1, k = 20, bs = "cc") + 
    s(WIND.2, k = 30, bs = "cc") + s(WIND.3, k = 30, bs = "cc") + 
    s(WINDSPD, k = 30) + factor(SKY) + s(MOON, k = 30) + factor(BIRD)

Parametric coefficients:
                    Estimate Std. Error t value Pr(>|t|)    
(Intercept)        2.880e+00  8.486e-02  33.940   <2e-16 ***
factor(SKY)cloudy -4.954e-05  9.023e-02  -0.001   0.9996    
factor(SKY)fog    -1.048e-02  1.310e-01  -0.080   0.9362    
factor(SKY)rain   -2.195e-01  1.541e-01  -1.425   0.1544    
factor(BIRD)2     -2.001e+00  1.916e-01 -10.444   <2e-16 ***
factor(BIRD)3     -2.603e+00  2.514e-01 -10.355   <2e-16 ***
factor(BIRD)4      2.065e-01  8.888e-02   2.324   0.0202 *  
factor(BIRD)5      1.228e+00  7.508e-02  16.356   <2e-16 ***
---
Signif. codes:  0 '***' 0.001 '**' 0.01 '*' 0.05 '.' 0.1 ' ' 1
\end{verbatim}
}
\end{frame}

\begin{frame}[fragile]{}

{\scriptsize
\begin{verbatim}
Approximate significance of smooth terms:
              edf Ref.df     F  p-value    
s(YEAR)    28.697 28.964 8.007  < 2e-16 ***
s(DAY)     17.522 20.852 5.485 2.00e-14 ***
s(NO.OBS)   3.534  4.265 9.785 4.15e-08 ***
s(WIND.0)  24.510 28.000 3.388 9.87e-12 ***
s(WIND.1)   8.080 18.000 2.438 1.30e-08 ***
s(WIND.2)   2.820 28.000 0.652 3.03e-05 ***
s(WIND.3)  23.924 28.000 1.585  0.00223 ** 
s(WINDSPD)  8.943 10.887 4.105 5.99e-06 ***
s(MOON)    26.886 28.474 2.840 8.02e-07 ***
---
Signif. codes:  0 '***' 0.001 '**' 0.01 '*' 0.05 '.' 0.1 ' ' 1

R-sq.(adj) =  0.103   Deviance explained =   67%
GCV = 43.889  Scale est. = 54.171    n = 2230
\end{verbatim}
}
\end{frame}


\begin{frame}[fragile]{Quasi-Poisson, base cyclique pour \texttt{WIND}, sans valeur aberrante}

{\scriptsize
\begin{verbatim}
> gam.check(GrA.quasifit.wo)
Method: GCV   Optimizer: outer newton
full convergence after 9 iterations.
Gradient range [-2.221881e-06,2.419194e-05]
(score 32.65511 & scale 44.22684).
Hessian positive definite, eigenvalue range [0.00213954,0.02999257].
Model rank =  163 / 163 

Basis dimension (k) checking results. Low p-value (k-index<1) may
indicate that k is too low, especially if edf is close to k'.
               k'    edf k-index p-value
s(YEAR)    19.000 18.826   0.568    0.00
s(DAY)     19.000 16.152   0.841    0.00
s(NO.OBS)   7.000  1.939   0.617    0.00
s(WIND.0)  18.000  8.973   0.858    0.00
s(WIND.1)  18.000 15.711   0.920    0.02
s(WIND.2)  18.000 15.715   0.887    0.00
s(WIND.3)  18.000 15.871   0.904    0.00
s(WINDSPD) 19.000 15.680   0.910    0.00
s(MOON)    19.000 17.462   0.968    0.55
\end{verbatim}
}
\end{frame}

\begin{frame}
\begin{center}
\includegraphics[width = 7.5cm]{figures/GrAquasicheckwo.pdf}  
\end{center}
\end{frame}


\begin{frame}[fragile]{Binomiale négative, base cyclique pour \texttt{WIND}, sans valeur aberrante}

{\scriptsize
\begin{verbatim}
> gam.check(GrAFitnegbin.wo)
Method: REML   Optimizer: outer newton
full convergence after 6 iterations.
Gradient range [-0.002532445,0.001774942]
(score 7125.691 & scale 1).
Hessian positive definite, eigenvalue range [0.002520513,740.6654].
Model rank =  153 / 153 

Basis dimension (k) checking results. Low p-value (k-index<1) may
indicate that k is too low, especially if edf is close to k'.
               k'    edf k-index p-value
s(YEAR)    19.000 14.481   0.356    0.00
s(DAY)     19.000  6.014   0.620    0.00
s(NO.OBS)   7.000  2.457   0.404    0.00
s(WIND.0)  18.000  3.852   0.703    0.00
s(WIND.1)  18.000  2.709   0.734    0.00
s(WIND.2)  18.000  2.182   0.689    0.00
s(WIND.3)  18.000  2.225   0.729    0.00
s(WINDSPD) 19.000  2.202   0.698    0.00
s(MOON)     9.000  1.007   0.778    0.34
\end{verbatim}
}
\end{frame}


\begin{frame}
\begin{center}
\includegraphics[width = 7.5cm]{figures/GrANBcheckwo.pdf}  
\end{center}
\end{frame}


\begin{frame}[fragile]{}

{\scriptsize
\begin{verbatim}
> summary(GrAFitnegbin.wo)

Family: Negative Binomial(0.444) 
Link function: log 

Formula:
NUMBER ~ s(YEAR, k = 20) + s(DAY, k = 20) + s(NO.OBS, k = 8) + 
    s(WIND.0, k = 20, bs = "cc") + s(WIND.1, k = 20, bs = "cc") + 
    s(WIND.2, k = 20, bs = "cc") + s(WIND.3, k = 20, bs = "cc") + 
    s(WINDSPD, k = 20) + factor(SKY) + s(MOON) + factor(BIRD)

Parametric coefficients:
                  Estimate Std. Error z value Pr(>|z|)    
(Intercept)        3.15046    0.08405  37.484   <2e-16 ***
factor(SKY)cloudy -0.06807    0.08731  -0.780   0.4356    
factor(SKY)fog    -0.19088    0.11097  -1.720   0.0854 .  
factor(SKY)rain   -0.12774    0.13333  -0.958   0.3380    
factor(BIRD)2     -2.42337    0.10868 -22.298   <2e-16 ***
factor(BIRD)3     -2.98799    0.11281 -26.487   <2e-16 ***
factor(BIRD)4     -1.03015    0.10348  -9.955   <2e-16 ***
factor(BIRD)5      1.52882    0.10154  15.057   <2e-16 ***
---
Signif. codes:  0 '***' 0.001 '**' 0.01 '*' 0.05 '.' 0.1 ' ' 1
\end{verbatim}
}
\end{frame}

\begin{frame}[fragile]{}

{\scriptsize
\begin{verbatim}

Approximate significance of smooth terms:
              edf Ref.df  Chi.sq  p-value    
s(YEAR)    14.481 16.591 188.028  < 2e-16 ***
s(DAY)      6.014  7.448 239.399  < 2e-16 ***
s(NO.OBS)   2.457  3.033  28.412 3.30e-06 ***
s(WIND.0)   3.852 18.000  75.064  < 2e-16 ***
s(WIND.1)   2.709 18.000  16.754 8.48e-05 ***
s(WIND.2)   2.182 18.000   9.214  0.00399 ** 
s(WIND.3)   2.225 18.000   9.895  0.00296 ** 
s(WINDSPD)  2.202  2.796  10.133  0.01661 *  
s(MOON)     1.007  1.013   0.067  0.79966    
---
Signif. codes:  0 '***' 0.001 '**' 0.01 '*' 0.05 '.' 0.1 ' ' 1

R-sq.(adj) =  -0.121   Deviance explained = 52.2%
-REML = 7125.7  Scale est. = 1         n = 2229
\end{verbatim}
}
\end{frame}


\begin{frame}
\begin{center}
\includegraphics[width = 8cm]{figures/GrAFitnegbinwo.pdf}  
\end{center}
\end{frame}

\begin{frame}[allowframebreaks]{Références du chapitre}
\tiny
\printbibliography[heading=none,section=\therefsection]
\end{frame}

\end{refsection}

\end{document}
